% Figure: the hydraulic press as force-multiplier. A small pumping piston and a
% large ram share one confined incompressible fluid at a common pressure p; the
% force ratio equals the area ratio while the travel ratio is its inverse, so
% work in equals work out. No em-dashes.
\begin{figure}[htbp]
\centering
\begin{tikzpicture}[scale=1.0]
% common fluid reservoir (the connecting body of oil)
\fill[engblue!15] (0.4,0.0) rectangle (8.6,1.2);
\draw[engblack, thick] (0.4,0.0) -- (8.6,0.0);
% small cylinder (pump) on the left
\draw[engblack, thick] (0.8,1.2) -- (0.8,3.4);
\draw[engblack, thick] (2.0,1.2) -- (2.0,3.4);
\fill[engblue!15] (0.8,1.2) rectangle (2.0,2.4);
% small piston
\fill[enggray!45] (0.8,2.4) rectangle (2.0,2.7);
\draw[engblack, very thick] (1.4,2.7) -- (1.4,3.6);
\draw[engamber, ultra thick, ->] (1.4,4.3) -- (1.4,3.7);
\node[engamber, font=\small, anchor=south] at (1.4,4.35) {$F_1$};
\node[engblack, font=\scriptsize, anchor=north] at (1.4,1.15) {area $A_1$};
% travel of the small piston (long stroke)
\draw[engred, thick, <->] (2.25,2.4) -- (2.25,3.5);
\node[engred, font=\scriptsize, anchor=west] at (2.3,2.95) {$\Delta_1$ (long)};
% large cylinder (ram) on the right
\draw[engblack, thick] (5.6,1.2) -- (5.6,3.4);
\draw[engblack, thick] (8.2,1.2) -- (8.2,3.4);
\fill[engblue!15] (5.6,1.2) rectangle (8.2,1.9);
% large piston (ram) with load
\fill[enggray!45] (5.6,1.9) rectangle (8.2,2.2);
\fill[englime!30] (6.0,2.2) rectangle (7.8,3.0);
\node[engblack, font=\small] at (6.9,2.6) {load};
\draw[englime, ultra thick, ->] (6.9,3.6) -- (6.9,4.2);
\node[englime, font=\small, anchor=south] at (6.9,4.25) {$F_2=pA_2$};
\node[engblack, font=\scriptsize, anchor=north] at (6.9,1.15) {area $A_2$};
% travel of the ram (short rise)
\draw[engred, thick, <->] (8.45,1.9) -- (8.45,2.2);
\node[engred, font=\scriptsize, anchor=west] at (8.5,2.05) {$\Delta_2$ (short)};
% common pressure label
\node[engblue, font=\small] at (4.0,0.6) {common pressure $p=F_1/A_1=F_2/A_2$};
\end{tikzpicture}
\caption[The hydraulic press]{The hydraulic press reads Pascal's principle as
a machine. One confined body of oil (blue) touches a small pumping piston of
area $A_1$ and a large ram of area $A_2$, and carries the same applied pressure
$p$ to both. The ram force is the pressure times its larger area, so
$F_2/F_1=A_2/A_1$, the force multiplication. The incompressible fluid conserves
volume, so the ram rises by the pumped volume divided by its area and its travel
$\Delta_2$ is smaller than the pump stroke $\Delta_1$ by the same area ratio.
Force out times distance out equals force in times distance in: the press buys
force with distance at a fixed rate, and stores no energy in the static ideal.}
\label{fig:vol03ch10:hydraulic-press}
\end{figure}
